% DELTA23 /root/five_level_six. An attempted main-gate detector that collapses.

\subsection*{D23.F11. The new symplectic covariance has only the old shear}
Let $t_h\in\widehat{\mathbb Z}$ be the full signed
three-strand shear of DELTA\_ON\_GT AD.1 for
$h\in\mathcal H=\ker\lambda\subset\widehat{GT}_0$.
Use the chain form and the representation $q$ of D23.T1, and
write $w=e_1+e_3$. Define
\[
 V(t)=\begin{pmatrix}
 1&t&0&t\\0&1&0&0\\0&0&1&t\\0&0&0&1
 \end{pmatrix}.
 \tag{D23.F11.1}
\]
Even without assuming the pentagon, the four Drinfeld
generator formulas, evaluated in this symplectic representation,
are exactly conjugation by $V(t_h)$.
For an actual full-GT element, this is the $V_h$ of D23.T1.
In particular
\[
 V_he_1=e_1,\quad V_he_3=e_3,\quad V_hw=w.
\]
The putative sign on the line $\widehat{\mathbb Z}w$ is
identically positive; it cannot detect $\Delta$.

\begin{proof}
The integral lattice is the orthogonal sum of the symplectic
planes
\[
 E=\langle e_1,e_2\rangle,\qquad
 E^\perp=\langle w,e_4\rangle,
\]
with each displayed ordered pair having pairing one.  The
matrix $V(t)$ is $\operatorname{diag}(U^t,U^t)$ in this
integral symplectic basis, so it preserves the chain form.

Put $t_i=q(\sigma_i)$ and
$z_i=(\sigma_1\cdots\sigma_{i-1})^i$.
The exact Drinfeld formula can be written
\[
 \sigma_1\mapsto\sigma_1,\qquad
 \sigma_i\mapsto f_h(z_i,\sigma_i^2)^{-1}
                   \sigma_i f_h(z_i,\sigma_i^2)\quad(i\geq2).
 \tag{D23.F11.2}
\]
See Nakamura, RIMS-1691, Section 5.4, equation (5.4.2),
printed p.~28.  This is equivalent to the usual formula
with $y_i=\sigma_{i-1}\cdots\sigma_1^2\cdots\sigma_{i-1}$:
$z_i=y_i z_{i-1}$, and $z_{i-1}$ commutes with both $y_i$
and $\sigma_i^2$.  Multiplying the first argument of a
commutator word by such a commuting central factor has no
effect, because its exponent sum in that argument is zero.

For $i=1$, $V(t)$ fixes $e_1$ and hence its transvection.
For $i=2$, the two arguments act on $E$ as $U^2,L^2$ and
fix $E^\perp$ pointwise.  AD.1 says precisely that the
resulting image of $t_2$ is its conjugate by $U^{t_h}$ on
$E$.  Thus it is conjugation by $V(t_h)$ on the full lattice.

For $i=3$, direct multiplication gives
\[
 q(z_3)=
 \begin{pmatrix}-1&0&2&0\\0&-1&0&0\\0&0&1&0\\0&0&0&1\end{pmatrix}.
\]
It preserves the plane $\langle e_1,w\rangle$, while
$t_3^2$ is identity on that plane.  Their restricted actions
commute.  Therefore the image of every commutator word,
in particular $f_h$, is identity on that plane and fixes
$e_3=w-e_1$.  The evaluated conjugate of $t_3$ is exactly
$t_3$, which is also its conjugate by $V(t_h)$.

For $i=4$, the integral identity
\[
 q(z_4)=(t_1t_2t_3)^4=t_w^2
\]
holds, where $t_w=I+ww^{\mathsf t}J$.
Both $t_w^2$ and $t_4^2$ act on $E^\perp$ as exactly
$U^2,L^2$, in the basis $(w,e_4)$, and fix $E$ pointwise.
The identical AD.1 calculation therefore gives the shear
$U^{t_h}$ on this second plane.  This is again conjugation
by $V(t_h)$.
All four formulas have now been checked.  For a full-GT
element they give the actual braid action, so uniqueness
in D23.T1 identifies $V_h$ with (D23.F11.1).  Its action
on $w$ is visibly identity.
\end{proof}

\paragraph{Scope of the failed detector.}
All the evaluated braid relations and this line sign are
automatic for gentle elements, including the extra-CM
element.  This does not assert that every separate lift
of the five-cycle word has trivial symplectic evaluation.
The representation does not descend to the five-point
sphere group, and any such lifted-pentagon assertion
requires its own calculation with its central factors retained.
